\section{Signed flow with moving comparison optima}
Let $\mP=\mA\mD^{-1}$, $\mK_0=(\mI+\mP)/2$, and
\[
 \beta_0=(1-\alpha)/(1+\theta)\leq\chi<\eta.
\]
The operator $\mK_0$ is nonnegative and column stochastic. In mass
coordinates write
\begin{equation}
 \begin{aligned}
 \vv_k&=\theta\mD^{1/2}\vz_k,
 &\va_k&=\theta\mD^{1/2}\vu_k,\\
 \vh_k^{\rm f}&=-\frac{\mD^{1/2}(\mQ-\mu\mI)
                                      (\vx_k-\vu_k)}{1+\theta},
 &\vs_k^{\rm f}&=\mD^{1/2}
                (\mQ\vu_k-\vb+\lambda_k\vomega).
 \end{aligned}
 \label{eq:revisit-flow-objects}
\end{equation}
The slack $\vs_k^{\rm f}$ is nonnegative. Direct substitution in the raw
update gives the exact fixed-comparator identity
\begin{equation}
 \theta\mD^{1/2}\vq_k-\va_k
       =\beta_0\mK_0(\vv_k-\va_k)+\vh_k^{\rm f}-\vs_k^{\rm f}.
 \label{eq:revisit-raw-flow}
\end{equation}
This is the same algebraic identity as in the source deterministic proof;
its comparator is now permitted to change after the step.

Set
\[
 \ve_k=(\vv_k-\va_k)_+,\qquad
 \mathcal I_k=\{i:v_{k+1,i}>a_{k+1,i}\}.
\]
Since $r_{k+1}\leq r_k$, monotonicity gives
$\va_{k+1}\geq\va_k$. A selected coordinate has positive new kinetic
state, so the orthant projection does not change its raw value.
Summing \cref{eq:revisit-raw-flow} on $\mathcal I_k$, and dropping the
favorable comparison shift $\va_{k+1}-\va_k$, proves
\begin{equation}
 \|\ve_{k+1}\|_1+\sum_{i\in\mathcal I_k}s_{k,i}^{\rm f}
   \leq\beta_0\|\ve_k\|_1
                   +\sum_{i\in\mathcal I_k}h_{k,i}^{\rm f}.
 \label{eq:revisit-selected-flow}
\end{equation}
Here positivity and column stochasticity bound the selected propagated
signed excess by its total positive mass. Retaining signed forcing on
the selected set is essential.

\begin{lemma}[The cooling weights preserve the telescope]
\label{lem:revisit-weighted-flow}
For every horizon $K$,
\begin{equation}
 \sum_{k<K}\sum_{i\in\mathcal I_k}
                          \frac{s_{k,i}^{\rm f}}{\lambda_k}
 \leq\sum_{k<K}\sum_{i\in\mathcal I_k}
                          \frac{h_{k,i}^{\rm f}}{\lambda_k}.
 \label{eq:revisit-weighted-flow}
\end{equation}
\end{lemma}
\begin{proof}
Divide \cref{eq:revisit-selected-flow} by $\lambda_k$ and sum.
The initial excess is zero because $\vz_0=\vu_0=\vzero$.
Every intermediate excess has coefficient
$1/\lambda_{k-1}-\beta_0/\lambda_k\geq0$, since
$\lambda_k\geq\eta\lambda_{k-1}\geq\beta_0\lambda_{k-1}$.
The terminal excess also has a nonnegative coefficient. Discard them.
\end{proof}

For each $k$ use the analytical core
$\mathcal C_k=\supp(\vx_{r_k/2}^*)$. It has volume at most $2/r_k$.
Off this core, the KKT slack satisfies
\begin{equation}
 i\notin\mathcal C_k\quad\Longrightarrow\quad
 (u_k)_i=(u_{k+1})_i=0,
 \qquad s_{k,i}^{\rm f}\geq\lambda_kd_i/2.
 \label{eq:revisit-core-margin}
\end{equation}
Indeed $r_{k+1}\geq\eta r_k\geq r_k/2$, so both comparison optima
vanish there. Comparing their zero row at $r_k$ with the optimum at
$r_k/2$ and using the nonpositive off-diagonals proves the slack margin.

\begin{theorem}[Cumulative kinetic degree work]
\label{thm:revisit-volume}
Every finite horizon of \cref{eq:revisit-recurrence} with the stated cooling
schedule satisfies
\begin{equation}
 \sum_{k<K}\vol(\supp\vz_{k+1})
                  \leq44\sum_{k<K}\frac1{r_k}.
 \label{eq:revisit-volume}
\end{equation}
This includes all excursions, first exposures, and repeated kinetic scans.
\end{theorem}
\begin{proof}
Put $S_K=\sum_{k<K}1/r_k$ and
\[
 V_{\rm out}=\sum_{k<K}\vol(\supp\vz_{k+1}\setminus\mathcal C_k),
 \qquad
 H_K=\sum_{k<K}
     \frac{\|\mD^{-1/2}\vh_k^{\rm f}\|_2^2}{\lambda_k^2}.
\]
Since $\mQ-\mu\mI$ commutes with $\mQ$ and its eigenvalues lie between
zero and those of $\mQ$, \cref{eq:revisit-optimum-response} gives
$H_K\leq10S_K$. Every positive kinetic coordinate outside its current
core is selected. By \cref{eq:revisit-weighted-flow,eq:revisit-core-margin}
and weighted Cauchy--Schwarz,
\[
 \tfrac12V_{\rm out}
 \leq\sum_{k<K}\sum_{i\in\mathcal I_k}h_{k,i}^{\rm f}/\lambda_k
 \leq\sqrt{(V_{\rm out}+2S_K)H_K}.
\]
The total selected degree volume is at most $V_{\rm out}+2S_K$,
regardless of how often a coordinate reappears. Set
$Y=V_{\rm out}+2S_K$. Young's inequality gives
\[
 Y\leq2S_K+2\sqrt{YH_K}
        \leq2S_K+Y/2+2H_K,
 \qquad Y\leq4S_K+4H_K\leq44S_K.
\]
The total kinetic volume is at most $Y$.
\end{proof}
